\section{Evaluation}

\subsection{Methodology}

We will evaluate the GPIS implementation using two key metrics: image quality and rendering performance.

\subsubsection{Image Quality Assessment}

We will use Mean Squared Error (MSE) to compare images rendered with a regular path tracer against those rendered with the GPIS object. MSE provides a pixel-by-pixel quantitative measure of the difference between two images, making it suitable for assessing whether the GPIS approach produces visually equivalent results to traditional mesh-based rendering.

MSE is calculated as:

\begin{equation}
\text{MSE} = \frac{1}{N} \sum (I_{\text{reference}} - I_{\text{gpis}})^2
\end{equation}

where $N$ is the total number of pixels, $I_{\text{reference}}$ represents pixel values from the standard path tracer, and $I_{\text{gpis}}$ represents pixel values from the GPIS renderer.

MSE is a viable metric for this evaluation because:

\begin{itemize}[noitemsep]
    \item It provides an objective, quantitative measure of image similarity
    \item It is widely used in rendering research for comparing ground truth images
    \item Lower MSE values directly correspond to closer visual matches
\end{itemize}

\subsubsection{Rendering Performance}
We will compare rendering times between the traditional mesh-based approach and the GPIS approach at three different sample rates: 8 samples per pixel (spp), 32 spp, and 128 spp. This will enable us to quantify the performance overhead of the GPIS approach relative to standard ray tracing. Based on the additional computational cost of Gaussian process evaluation and correlated sampling, I anticipate the GPIS approach will be 5-10x slower than traditional mesh rendering.
